% !TEX TS-program = XeLaTeX
% use the following command:
% all document files must be coded in UTF-8
\documentclass[spanish]{textolivre}
% build HTML with: make4ht -e build.lua -c textolivre.cfg -x -u article "fn-in,svg,pic-align"
\usepackage{array}

\journalname{Texto Livre}
\thevolume{19}
%\thenumber{1} % old template
\theyear{2026}
\receiveddate{\DTMdisplaydate{2025}{7}{3}{-1}} % YYYY MM DD
\accepteddate{\DTMdisplaydate{2025}{12}{4}{-1}}
\publisheddate{\DTMdisplaydate{2026}{5}{14}{-1}}
\corrauthor{Irene Casanova-Mata}
\articledoi{10.1590/1983-3652.2026.60109}
%\articleid{NNNN} % if the article ID is not the last 5 numbers of its DOI, provide it using \articleid{} commmand 
% list of available sesscions in the journal: articles, dossier, reports, essays, reviews, interviews, editorial
\articlesessionname{articles}
\runningauthor{Casanova-Mata} 
%\editorname{Leonardo Araújo} % old template
\sectioneditorname{Hugo Heredia Ponce~\orcid{0000-0003-3657-1369}}
\layouteditorname{Saula Cecília~\orcid{0009-0006-3069-8480}}

\title{Impacto de la gamificación en el aula de inglés en 2.º de Educación Primaria desde una perspectiva de género}
% if there is a third language title, add here:
\othertitle{Impacto da gamificação na sala de aula de inglês no 2.º ano do Ensino Primário a partir de uma perspectiva de gênero}
\othertitle{Impact of gamification in the English classroom in 2nd Grade of Primary Education from a gender perspective}

\author[1]{Irene Casanova-Mata~\orcid{0009-0005-0445-2260}\thanks{Email: \href{mailto:profesor.icasanova@uclm.es}{profesor.icasanova@uclm.es}}}
\affil[1]{Universidad de Castilla-La Mancha (UCLM), Facultad de Educación Lorenzo Luzuriaga, Ciudad Real, España.}

\addbibresource{article.bib}
% use biber instead of bibtex
% $ biber article

% used to create dummy text for the template file
\definecolor{dark-gray}{gray}{0.35} % color used to display dummy texts
\usepackage{lipsum}
\SetLipsumParListSurrounders{\colorlet{oldcolor}{.}\color{dark-gray}}{\color{oldcolor}}

% used here only to provide the XeLaTeX and BibTeX logos
\usepackage{hologo}

% if you use multirows in a table, include the multirow package
\usepackage{multirow}

% provides sidewaysfigure environment
\usepackage{rotating}

% CUSTOM EPIGRAPH - BEGIN 
%%% https://tex.stackexchange.com/questions/193178/specific-epigraph-style
\usepackage{epigraph}
\renewcommand\textflush{flushright}
\makeatletter
\newlength\epitextskip
\pretocmd{\@epitext}{\em}{}{}
\apptocmd{\@epitext}{\em}{}{}
\patchcmd{\epigraph}{\@epitext{#1}\\}{\@epitext{#1}\\[\epitextskip]}{}{}
\makeatother
\setlength\epigraphrule{0pt}
\setlength\epitextskip{0.5ex}
\setlength\epigraphwidth{.7\textwidth}
% CUSTOM EPIGRAPH - END

% to use IPA symbols in unicode add
%\usepackage{fontspec}
%\newfontfamily\ipafont{CMU Serif}
%\newcommand{\ipa}[1]{{\ipafont #1}}
% and in the text you may use the \ipa{...} command passing the symbols in unicode

% LANGUAGE - BEGIN
% ARABIC
% for languages that use special fonts, you must provide the typeface that will be used
% \setotherlanguage{arabic}
% \newfontfamily\arabicfont[Script=Arabic]{Amiri}
% \newfontfamily\arabicfontsf[Script=Arabic]{Amiri}
% \newfontfamily\arabicfonttt[Script=Arabic]{Amiri}
%
% in the article, to add arabic text use: \textlang{arabic}{ ... }
%
% RUSSIAN
% for russian text we also need to define fonts with support for Cyrillic script
% \usepackage{fontspec}
% \setotherlanguage{russian}
% \newfontfamily\cyrillicfont{Times New Roman}
% \newfontfamily\cyrillicfontsf{Times New Roman}[Script=Cyrillic]
% \newfontfamily\cyrillicfonttt{Times New Roman}[Script=Cyrillic]
%
% in the text use \begin{russian} ... \end{russian}
% LANGUAGE - END

% EMOJIS - BEGIN
% to use emoticons in your manuscript
% https://stackoverflow.com/questions/190145/how-to-insert-emoticons-in-latex/57076064
% using font Symbola, which has full support
% the font may be downloaded at:
% https://dn-works.com/ufas/
% add to preamble:
% \newfontfamily\Symbola{Symbola}
% in the text use:
% {\Symbola }
% EMOJIS - END

% LABEL REFERENCE TO DESCRIPTIVE LIST - BEGIN
% reference itens in a descriptive list using their labels instead of numbers
% insert the code below in the preambule:
%\makeatletter
%\let\orgdescriptionlabel\descriptionlabel
%\renewcommand*{\descriptionlabel}[1]{%
%  \let\orglabel\label
%  \let\label\@gobble
%  \phantomsection
%  \edef\@currentlabel{#1\unskip}%
%  \let\label\orglabel
%  \orgdescriptionlabel{#1}%
%}
%\makeatother
%
% in your document, use as illustraded here:
%\begin{description}
%  \item[first\label{itm1}] this is only an example;
%  % ...  add more items
%\end{description}
% LABEL REFERENCE TO DESCRIPTIVE LIST - END


% add line numbers for submission
%\usepackage{lineno}
%\linenumbers

\begin{document}
\maketitle

\begin{polyabstract}
\begin{abstract}
El objetivo de este estudio es evaluar el impacto de una metodología gamificada con el apoyo de herramientas digitales sobre el rendimiento en lengua inglesa de 96 estudiantes de 2.º de Educación Primaria (edad media: 8,25 años; DE = 0,44), con especial atención a las diferencias de género. Para ello, se adoptó un enfoque cuantitativo, con un diseño cuasi-experimental de tipo pretest-postest, empleando el \textit{Cuestionario de Competencias Lingüísticas Específicas} como instrumento de medida. La muestra se dividió en dos grupos equivalentes en cuanto a género: un grupo de control (GC) sometido a una metodología tradicional y un grupo experimental (GE) expuesto a la gamificación. En la muestra global no se observaron diferencias significativas de género durante el pretest (p = 0,769), aunque las niñas mostraron puntuaciones ligeramente superiores a las de los niños. Tampoco aparecieron diferencias significativas entre el GE y el GC en el pretest (p > 0,05). Sin embargo, en el postest los niños del GE obtuvieron medias significativamente superiores a sus pares del GC (p < 0,001) tras la intervención, y las niñas aumentaron su rendimiento con la gamificación, pero no significativamente (p = 0,319). Además, en el GC las niñas superaron significativamente a los niños en el postest, mientras que en el GE no se observaron diferencias significativas basadas en el género tras la intervención gamificada, lo que sugiere que este enfoque nivela brechas de género y puede ayudar a promover la equidad en el aprendizaje del inglés como lengua extranjera.
\keywords{Gamificación\sep Enseñanza del inglés\sep Equidad de género\sep Educación Primaria}
\end{abstract}

\begin{portuguese}
\begin{abstract}
O objetivo deste estudo é avaliar o impacto de uma metodologia gamificada, com o apoio de ferramentas digitais, no desempenho em língua inglesa de 96 alunos do 2.º ano do Ensino Primário (idade média: 8,25 anos; DP = 0,44), com especial atenção às diferenças de género. Para isso, adotou-se uma abordagem quantitativa, com um desenho quase-experimental do tipo pré-teste/pós-teste, utilizando o \textit{Questionário de Competências Linguísticas Específicas} como instrumento de medida. A amostra foi dividida em dois grupos equivalentes em termos de gênero: um grupo controle (GC), submetido a uma metodologia tradicional, e um grupo experimental (GE), exposto à gamificação. Na amostra geral, não foram observadas diferenças significativas de gênero no pré-teste (p = 0,769), embora as meninas tenham apresentado pontuações ligeiramente superiores às dos meninos. Também não foram encontradas diferenças significativas entre o GE e o GC no pré-teste (p > 0,05). No entanto, no pós-teste, os meninos do GE obtiveram médias significativamente superiores às de seus pares do GC (p < 0,001) após a intervenção, e as meninas aumentaram seu desempenho com a gamificação, embora sem significância estatística (p = 0,319). Além disso, no GC, as meninas superaram significativamente os meninos no pós-teste, enquanto no GE não foram observadas diferenças significativas baseadas no gênero após a intervenção gamificada, o que sugere que essa abordagem contribui para nivelar as disparidades de gênero e pode favorecer a promoção da equidade na aprendizagem do inglês como língua estrangeira.
\keywords{Gamificação\sep Ensino de inglês\sep Equidade de gênero\sep Ensino Fundamental}
\end{abstract}
\end{portuguese}

\begin{english}
\begin{abstract}
The objective of this study is to evaluate the impact of a gamified methodology, supported by digital tools, on the English language performance of 96 second-grade primary school students (mean age: 8.25 years; SD=0.44), with a specific focus on gender differences. To this end, a quantitative approach was adopted using a quasi-experimental pretest-posttest design, employing the Specific Linguistic Competencies Questionnaire as the measurement instrument. The sample was divided into two gender-equivalent groups: a control group (CG) subjected to traditional methodology and an experimental group (EG) exposed to gamification. In the global sample, no significant gender differences were observed during the pretest (p=0.769), although girls showed slightly higher scores than boys. Similarly, no significant differences appeared between the EG and the CG in the pretest (p>0.05). However, in the posttest, boys in the EG obtained significantly higher means than their peers in the CG (p<0.001) following the intervention. While girls increased their performance with gamification, the improvement was not statistically significant (p=0.319). Furthermore, in the CG, girls significantly outperformed boys in the posttest, whereas no significant gender-based differences were observed in the EG after the gamified intervention. This suggests that such an approach bridges gender gaps and may help promote equity in English as a Foreign Language (EFL) learning.
\keywords{Gamification\sep English Language Teaching (ELT)\sep Gender Equity\sep Primary Education}
\end{abstract}
\end{english}
% if there is another abstract, insert it here using the same scheme
\end{polyabstract}

\section{Introducción}\label{sec-intro}
El aprendizaje de lenguas extranjeras, especialmente del inglés, se ha consolidado como una prioridad educativa en contextos cada vez más globalizados y multilingües debido a necesidades profesionales, académicas y socioculturales derivadas de la internacionalización del conocimiento \cite{arrieta2023}. En el caso específico del sistema educativo español, la Ley Orgánica 3/2020, de 29 de diciembre, por la que se modifica la Ley Orgánica 2/2006, de 3 de mayo \cite{espana2020}, subraya la importancia de una formación plurilingüe, alineada con una ciudadanía ``moderna, más abierta, menos rígida, multilingüe y cosmopolita'' \cite[p. 6]{espana2020}. Esta tendencia ha ido acompañada de importantes cambios en la enseñanza de lenguas, como la adopción del Marco Común Europeo de Referencia para las Lenguas (MCER) y la incorporación de enfoques comunicativos orientados a tareas motivadoras \cite{cambridge2011}. Sin embargo, dichos enfoques no siempre concretan las metodologías más adecuadas para garantizar aprendizajes significativos e inclusivos.

En este contexto, las metodologías activas, apoyadas por las herramientas tecnológicas, se presentan como alternativas relevantes para la enseñanza de segundas lenguas \cite{silva2017, alvarez2021}. Entre ellas, la gamificación destaca como una estrategia pedagógica que integra elementos del juego en contextos educativos con el fin de aumentar la motivación, el compromiso y el rendimiento académico \cite{zambrano2020}. Diversos estudios evidencian su eficacia en el aprendizaje del inglés \cite{li2022, lampropoulos2022}, así como sus efectos positivos en el clima de aula y en la actitud del alumnado hacia la lengua extranjera \cite{sanchez2019, mora2020}. Para que la gamificación tenga éxito, es necesario desarrollar una narrativa coherente y contextualizada que fomente la implicación activa del alumnado \cite{quintero2018}.

No obstante, la investigación sobre gamificación educativa aún muestra una laguna significativa en relación con su impacto desde una perspectiva de género. Si bien algunos estudios consideran la gamificación una herramienta pedagógica neutral \cite{zahedi2021}, hay investigaciones que demuestran que las mujeres tienden a ser mejores aprendices de una segunda lengua que los hombres \cite{barrios2015, glowka2014}. Este hecho plantea interrogantes sobre si la gamificación puede contribuir a reducir posibles brechas de género en el aprendizaje lingüístico. 

A partir de estas premisas, el objetivo principal de este estudio es analizar la incidencia de la gamificación apoyada por herramientas tecnológicas en el rendimiento académico en lengua inglesa de estudiantes de Educación Primaria (EP), prestando especial atención a la variable de género, definida en función del sexo biológico. Para ello, se plantean las siguientes preguntas de investigación (PI):

\begin{quote}
    PI1: ¿Existen diferencias significativas en el rendimiento académico en lengua inglesa entre niños y niñas antes y después de la implementación de diferentes metodologías de enseñanza?
\\
    PI2: ¿Se observan diferencias significativas de género en función del tipo de metodología empleada (tradicional vs. gamificada)?
\\
PI3: ¿Qué impacto tiene la metodología aplicada (tradicional vs. gamificada) en el rendimiento académico de los estudiantes varones?
\\
PI4: ¿Qué impacto tiene la metodología aplicada (tradicional vs. gamificada) en el rendimiento académico de las estudiantes mujeres?
\end{quote}

La estructura del presente artículo se organiza del siguiente modo: en primer lugar, se desarrolla el marco teórico sobre gamificación y rendimiento lingüístico desde una perspectiva pedagógica y de género. A continuación, se describe el diseño metodológico. Posteriormente, se presentan los resultados y su discusión en relación con la literatura previa. Finalmente, se exponen las conclusiones, las implicaciones pedagógicas y las líneas futuras de investigación derivadas del estudio.

\section{Marco teórico}

\subsection{La enseñanza del inglés en contextos educativos actuales}
En las últimas décadas, el inglés se ha consolidado como la lengua franca por excelencia en una amplia gama de dominios internacionales, como el ámbito empresarial, la ciencia y la tecnología, así como en las relaciones diplomáticas y los entornos académicos \cite{arrieta2023}. Su posición privilegiada como lengua vehicular de la globalización ha contribuido a que el dominio funcional del inglés sea considerado una competencia transversal clave para el siglo XXI. En este contexto, el acceso temprano y equitativo a su enseñanza cobra una relevancia estratégica en los sistemas educativos contemporáneos.

En el caso de España, la reforma educativa articulada a través de la \textcite{espana2020} refuerza el enfoque plurilingüe como uno de los ejes vertebradores del currículo, al establecer la necesidad de formar ciudadanos abiertos, críticos y cosmopolitas. Esta orientación normativa impulsa la implementación de programas de lengua extranjera desde las etapas iniciales y aboga por una educación más flexible, inclusiva y centrada en el desarrollo de competencias comunicativas. En este marco, la enseñanza del inglés no se concibe únicamente como una materia instrumental, sino como un vehículo para la participación activa en entornos globales.

Paralelamente, el MCER, promovido por el Consejo de Europa, se ha convertido en el referente metodológico y evaluativo predominante en el diseño de itinerarios de aprendizaje lingüístico. Este marco proporciona una escala común para describir los niveles de competencia comunicativa e impulsa prácticas docentes orientadas al uso real del idioma en situaciones significativas \cite{gomez2020}. No obstante, aunque el MCER establece descriptores de logro y principios evaluativos, no prescribe metodologías específicas, lo que otorga a los agentes educativos autonomía para incorporar enfoques didácticos innovadores.

Entre esos enfoques, los modelos de enseñanza por tareas han ido ganando terreno en la didáctica del inglés como lengua extranjera, favoreciendo una aproximación más funcional y contextualizada del idioma \cite{cambridge2011}. Estas metodologías promueven la interacción auténtica, la negociación de significado y la resolución de problemas comunicativos, lo cual contribuye a desarrollar tanto la fluidez como la precisión lingüística. Sin embargo, este enfoque basado en el aprendizaje activo exige también que las propuestas sean significativas desde el punto de vista afectivo y motivacional para el alumnado, especialmente en etapas como la EP, donde el componente lúdico desempeña un papel esencial en la disposición al aprendizaje.

Por esta razón, en el marco de un sistema educativo que busca adaptarse a los retos de la sociedad global del conocimiento, resulta prioritario investigar e implementar estrategias didácticas que no solo aseguren el cumplimiento de estándares lingüísticos, sino que también potencien el compromiso, la motivación y la equidad en el aprendizaje del inglés como lengua extranjera. La gamificación, concebida como metodología activa emergente, ofrece un camino prometedor para lograr este equilibrio pedagógico.

\subsection{Metodologías activas y gamificación en el aula de lengua extranjera}
En el contexto actual, las metodologías activas se consolidan como enfoques centrados en el estudiante, en los que el aprendizaje se construye mediante la exploración, la reflexión crítica y la interacción significativa con el entorno \cite{silva2017}. Estas metodologías priorizan la participación activa del alumnado en tareas contextualizadas y con propósito, que estimulen tanto el desarrollo cognitivo como socioemocional \cite{londono2020}. En el ámbito de la enseñanza de lenguas extranjeras, el uso de metodologías activas implica abandonar paradigmas transmisivos centrados en la memorización de estructuras gramaticales, para dar paso a procesos de negociación de significado, producción auténtica del discurso y colaboración entre iguales \cite{alvarez2021}.

Bajo esta premisa, la gamificación emerge como una metodología de aprendizaje activo y participativo que recurre al diseño de experiencias inspiradas en el juego para implicar al estudiantado de manera significativa \cite{contreras2016}. Entendida como la incorporación de elementos característicos del juego – como mecánicas, dinámicas y estéticas – en contextos no lúdicos \cite{zambrano2020}, la gamificación se alinea con los postulados del aprendizaje experiencial y motivacional. \textcite{parra2018} la describen como el uso estratégico de recursos lúdicos orientados al logro de objetivos académicos en contextos reales, capaces de generar experiencias de aprendizaje significativas. Asimismo, estos recursos, junto con herramientas digitales, potencian la motivación intrínseca, promueven la interacción social y contribuyen a la mejora del desempeño oral del alumnado \cite{wong2021}. Estas aproximaciones coinciden en señalar que la gamificación no se limita a la introducción superficial de recompensas, sino que requiere un diseño pedagógico intencional basado en retos progresivos, retroalimentación constante y coherencia didáctica, favoreciendo procesos cognitivos complejos y funciones ejecutivas como la memoria de trabajo y la flexibilidad cognitiva \cite{bedoya2022}.

Desde una perspectiva aplicada, su aplicación en el aula de lengua extranjera permite transformar actividades convencionales en desafíos significativos que incluyen recompensas, niveles, misiones, sistemas de puntuación y retroalimentación inmediata. Más allá de su componente instrumental, su eficacia se vincula con la activación de procesos emocionales positivos, el fomento de la autonomía y la generación de sentido de pertenencia \cite{garciamoreno2019, cordero2020}. Además, elementos como avatares, insignias, clasificaciones colaborativas o narrativas interactivas tecnológicas facilitan la inmersión del alumnado, incrementando la atención sostenida y la resiliencia ante el error \cite{iqbal2023}.

La evidencia empírica apoya firmemente el valor pedagógico de la gamificación en el aprendizaje de lenguas extranjeras, posicionándola como un catalizador del aprendizaje significativo y de la implicación activa del alumnado. Diversos estudios han demostrado que su implementación puede mejorar el rendimiento académico del alumnado \cite{sanchez2019, fiestas2023} y aumentar su motivación intrínseca \cite{lampropoulos2022} mediante la participación en experiencias de aprendizaje estructuradas y desafiantes. Del mismo modo, la literatura destaca su impacto positivo en el clima de aula, al favorecer dinámicas cooperativas y el desarrollo del desempeño oral mediante la interacción interpersonal \cite{wong2021}, así como su contribución a la regulación emocional y a la reducción de la ansiedad comunicativa en contextos de uso de la lengua extranjera \cite{mora2020}. Estos efectos se potencian cuando la gamificación está sustentada por apoyo tecnológico, una narrativa coherente, un diseño instruccional adaptado al nivel de competencia del alumnado y una secuenciación progresiva de los retos \cite{quintero2018}. En consecuencia, la gamificación no debe concebirse como una estrategia meramente decorativa o superficial, sino como una metodología integral con capacidad para transformar el aula de lengua extranjera en un espacio estimulante, emocionalmente seguro y pedagógicamente riguroso, favoreciendo el desarrollo integral del estudiante \cite{garcia2021}.

\subsection{Gamificación y equidad de género en el aprendizaje del inglés}
El vínculo entre innovación metodológica y equidad de género constituye una línea de investigación cada vez más relevante en el ámbito educativo, especialmente en lo que concierne al aprendizaje de lenguas extranjeras \cite{garciacalvo2022, nieto2019, nieto-casanova2026}. Con frecuencia, se evidencia una tendencia favorable al rendimiento lingüístico femenino, asociando este fenómeno a factores como la motivación intrínseca, la mayor autorregulación académica o una actitud positiva hacia el aprendizaje de lenguas \cite{porvoyer2014, barrios2015}. Diversos estudios han señalado que, a diferencia de las niñas, algunos estudiantes varones suelen partir de niveles más reducidos de implicación en el aprendizaje de lenguas extranjeras, expresando una inclinación más favorable hacia metodologías que incorporan elementos prácticos o de carácter lúdico para fomentar su implicación \cite{fernandez2014, lee2017}. Esta asimetría ha contribuido a reforzar ciertas narrativas esencialistas sobre las capacidades lingüísticas por género, que deben ser analizadas con cautela y superadas desde una perspectiva pedagógica basada en la equidad.

En este contexto, la gamificación emerge como una estrategia metodológica con potencial para compensar dichas disparidades percibidas. Al basarse en dinámicas centradas en la participación activa, el refuerzo positivo y la retroalimentación inmediata, esta metodología permite activar distintos perfiles motivacionales y promover un acceso más equitativo a los aprendizajes \cite{zahedi2021}. Más aún, el diseño de experiencias gamificadas contextualizadas – cuando se apoya en narrativas atractivas, progresión por niveles y metas significativas – puede favorecer una implicación emocional sostenida, especialmente en alumnado que requiere estímulos extrínsecos para mantener el compromiso \cite{garciamoreno2019}. Investigaciones como la de \textcite{liu2024} apuntan a que este enfoque puede conectar de forma particularmente eficaz con estudiantes varones, quienes a menudo encuentran en el juego una vía más accesible para activar sus recursos comunicativos.

Por otro lado, estudios como el de \textcite{fernandezbarrionuevo2022} señalan que las niñas tienden a mantener una actitud favorable y constante hacia el aprendizaje del inglés, independientemente del tipo de metodología empleada. Este hallazgo podría sugerir la existencia de un posible ``efecto de techo'' metodológico, en el que la motivación inicial elevada del alumnado femenino condiciona la posibilidad de mejoras significativas atribuibles exclusivamente a la gamificación. No obstante, investigaciones recientes subrayan que el verdadero potencial transformador de la gamificación no reside solo en su impacto inmediato sobre el rendimiento, sino en su capacidad para crear entornos de aprendizaje emocionalmente seguros, colaborativos y libres de sesgos \cite{mellado2024}.

Así, resulta pertinente destacar que la equidad educativa no se limita al acceso igualitario a los contenidos curriculares, sino que implica también la adecuación de las prácticas metodológicas a las características, necesidades e intereses del alumnado. Esta idea se alinea con el principio de personalización del aprendizaje que promueven tanto el MCER como las últimas reformas curriculares en España, en particular aquellas que defienden un enfoque competencial, inclusivo y motivador. Siguiendo estas líneas, la presente investigación busca aportar evidencia empírica sobre el papel de la gamificación como herramienta para reducir – o al menos equilibrar – posibles brechas de rendimiento asociadas al género en el aprendizaje del inglés en EP, contribuyendo al debate sobre cómo combinar innovación didáctica y justicia pedagógica en contextos escolares contemporáneos.

\section{Metodología}
La presente investigación adopta un enfoque cuantitativo, con un diseño cuasi-experimental de tipo pretest-postest mediante un cuestionario \textit{ad hoc} validado. Las variables independientes corresponden a la implementación de la gamificación en el aula y al género del alumnado, limitado a la variable género, entendida como sexo biológico (niña, niño). La variable dependiente se define a partir del rendimiento académico en lengua inglesa, evaluado a través de competencias lingüísticas específicas: comprensión oral y escrita, así como expresión oral y escrita.

\subsection{Muestra}
El presente estudio se llevó a cabo en cuatro centros de EP ubicados en la provincia de Ciudad Real, caracterizados por un nivel socioeconómico medio-alto, según consta en sus respectivos Proyectos Educativos. La muestra seleccionada estuvo compuesta por alumnado de 2.º curso de EP (n = 96; 50\% niños; 50\% niñas), con una edad media de 8,25 años (DE = 0,44). Como criterio principal de inclusión, se exigió la firma del consentimiento informado por parte de los representantes legales del alumnado participante.

La investigación se basó en un muestreo no probabilístico por conveniencia, justificado por razones de accesibilidad y factibilidad, dado el contacto previo y la disponibilidad de los centros para la intervención. El grupo experimental (n = 48; 50\% niños y 50\% niñas), con una edad media de 8,33 años (DE = 0,49), recibió una intervención didáctica basada en gamificación. Por su parte, el grupo de control (n = 48; 50\% niños y 50\% niñas), con una edad media de 8,17 años (DE = 0,49), siguió una metodología tradicional centrada en actividades y ejercicios del libro de texto. 

\subsection{Instrumento}
El instrumento utilizado en esta investigación fue el \textit{Cuestionario de Competencias Lingüísticas Específicas (CCLE)} para evaluar el rendimiento académico del alumnado en lengua inglesa. Su elaboración se fundamentó en los contenidos curriculares establecidos en la programación didáctica, e incluyó aspectos léxicos y gramaticales como el vocabulario relacionado con animales de granja y las estructuras \textit{there is / there are} y \textit{It's got / It hasn't got}.

El cuestionario consta de 117 ítems distribuidos en cuatro secciones correspondientes a las competencias lingüísticas evaluadas: comprensión escrita (26 ítems), expresión escrita (58 ítems), comprensión oral (16 ítems) y expresión oral (17 ítems). La Tabla \ref{tab-1} muestra las actividades utilizadas en cada competencia lingüística y los ítems de respuesta correspondientes. 

%--- codigo da tabela 1 ---%
\begin{longtable}{>{\raggedright\arraybackslash}p{1.9cm} >{\raggedright\arraybackslash}p{1.2cm} >{\raggedright\arraybackslash}p{1cm} p{9cm}}
\caption{Descripción de las actividades del \textit{CCLE}.}\label{tab-1} \\

\toprule
Competencia lingüística & Actividad & Ítems & Descripción \\
\midrule
\endfirsthead

\endhead

\endfoot

\bottomrule
\source{Elaboración propia.}
\endlastfoot

Comprensión oral (listening) & CO1 & 8 & Discriminación auditiva para actividad de elección múltiple con soporte visual: El alumnado escucha breves descripciones y selecciona la opción correcta entre tres imágenes \\
& CO2 & 1 & Identificación auditiva global de asociación semántica con opción múltiple extendida: El alumnado debe marcar una de las seis imágenes a partir de un audio descriptivo \\
& CO3 & 7 & Ordenación secuencial auditiva: El alumnado debe organizar una serie de animales según el orden en que aparecen en el audio \\

\addlinespace
Comprensión escrita (reading) & CE1 & 6 & Lectura comprensiva y asociación visual: El alumnado lee descripciones breves que integran la estructura gramatical \textit{there is / there are} para colorear una imagen \\
& CE2 & 6 & Lectura comprensiva con discriminación semántica y morfosintáctica: Mediante el uso de \textit{it’s got / it hasn’t got}, el alumnado analiza la presencia o ausencia de partes del cuerpo de animales \\
& CE3 & 6 & Lectura comprensiva con producción gráfica: A partir de un texto que indica qué animales hay en una granja, el alumnado realiza un dibujo representando aquellos animales mencionados \\
& CE4 & 8 & Lectura selectiva con opción visual dual: El alumnado lee una lista de animales presentes en una granja y, con base en esa lectura, debe seleccionar entre dos imágenes (A o B) la que se ajusta a la información proporcionada en el texto \\

\addlinespace

Expresión oral (speaking) &  &  & Descripción oral libre a partir de una imagen de animales en una granja, evaluada mediante cuatro indicadores analíticos \\
& EO1 & 8 & Fluidez léxica: Se asigna una puntuación entre 0 y 8 puntos en función del número de animales identificados y nombrados adecuadamente \\
& EO2 & 5 & Precisión léxica: Se otorgan hasta 5 puntos según el número de elementos anatómicos correctamente identificados en la producción oral \\
& EO3 & 2 & Corrección gramatical \textit{(there is / there are)}: Se valora con 2 puntos si se utilizan correctamente y de forma reiterada durante la descripción; 1 punto si su uso es esporádico; 0 puntos si no se emplean \\
& EO4 & 2 & Corrección gramatical \textit{(it’s got / it hasn’t got)}: Se califica también con un máximo de 2 puntos bajo los mismos criterios anteriores \\

\addlinespace

Expresión escrita (writing) & EE1 & 12 & Reconocimiento ortográfico con apoyo visual: A partir de una imagen de un animal, el alumnado debe reorganizar letras desordenadas para formar correctamente su nombre \\
& EE2 & 15 & Escritura guiada con concordancia gramatical: El alumnado observa una imagen de un animal y redacta una oración indicando lo que tiene y lo que no tiene, utilizando \textit{it’s got...} o \textit{it hasn’t got...} \\
& EE3 & 24 & Descripción cuantitativa con estímulo visual: El alumnado debe redactar oraciones describiendo cuántos animales hay en cada imagen, empleando \textit{there is / there are} y cuantificadores adecuados \\
& EE4 & 7 & Escritura deductiva basada en inferencia textual: A partir de una descripción escrita de las características de un animal, el alumnado deduce de qué animal se trata y escribe su nombre \\

\end{longtable}

El rendimiento académico se cuantificó a través de la suma total de las puntuaciones obtenidas en cada sección del cuestionario, estableciendo un rango de puntuación total entre 0 y 117 puntos.

Para garantizar la validez del instrumento, este fue sometido a un proceso de evaluación por juicio de expertos que incluyó a 6 profesores universitarios expertos en lengua inglesa, así como a 5 maestros de EP, especialistas en lengua extranjera (inglés). Los expertos evaluaron el instrumento atendiendo a los criterios de adecuación, pertinencia, claridad y coherencia con los objetivos del estudio, siguiendo las rúbricas del \textit{Formato de Validación por Expertos. Guía para validar instrumentos de investigación} \cite{direccion2018}.

Con el fin de determinar la validez de contenido del instrumento mediante juicio de expertos, se calculó el coeficiente de concordancia de Kendall (W), obteniéndose un valor de W = 0,650, lo que indica un alto grado de concordancia interjueces, estadísticamente significativo (p < 0,001), de acuerdo con la escala de interpretación propuesta por \textcite{escobar2008}. A partir de las aportaciones de los expertos, se realizaron ajustes en la redacción de instrucciones, incorporando un lenguaje más sencillo para facilitar la comprensión del instrumento por parte del alumnado de EP. Este proceso aseguró la adecuación pedagógica y lingüística del instrumento a la muestra y a los objetivos del estudio.

\subsection{Procedimiento}
La presente investigación se desarrolló durante el mes de abril de 2023 en fases secuenciales. En primer lugar, se celebró una reunión con el Equipo Directivo de los centros educativos seleccionados con el fin de seleccionar la muestra participante. Tras la elaboración del CCLE y el diseño de la experiencia didáctica gamificada, se informó del proyecto a los representantes legales del alumnado. Una vez obtenida la autorización correspondiente mediante consentimiento, se procedió a la aplicación del pretest (CCLE) a los dos grupos implicados (experimental y control), con carácter previo a la implementación de las intervenciones didácticas: tradicional y gamificada. 

La intervención gamificada fue concebida como una propuesta didáctica estructurada en 8 sesiones, cada una de 45 minutos, enmarcada en una narrativa digital interactiva que situaba al alumnado en una misión cooperativa. El eje de dicha narrativa giraba en torno a escenarios digitales interrelacionados, a los que los equipos solo podían acceder si lograban superar los desafíos lingüísticos propuestos en cada nivel. La dinámica se articuló de forma cooperativa: el alumnado trabajaba en pequeños grupos que debían completar tareas académicas, demostrar actitudes de ayuda mutua, expresarse en lengua inglesa y respetar las normas de la dinámica para acumular puntos y avanzar.

Este sistema de puntos, gestionado mediante la plataforma digital ClassDojo, permitía un seguimiento en tiempo real del progreso de cada grupo, así como reforzar conductas positivas a través de la entrega de insignias, premios simbólicos y retroalimentación constante. Los puntos podían obtenerse mediante la resolución exitosa de actividades lingüísticas, el cumplimiento de tareas complementarias y la demostración de actitudes colaborativas. Cada dos niveles superados, los equipos accedían a recompensas previamente anunciadas, generando así expectativas motivadoras que mantenían el interés sostenido a lo largo de todo el proceso. 

Cada sesión se organizó en torno a una secuencia pedagógica constante: presentación del reto diario, desarrollo de actividades lingüísticas interactivas vinculadas a los animales de granja, puntuación y reconocimiento del rendimiento grupal e individual, y formulación de tareas voluntarias para profundizar en los contenidos.

Una característica distintiva de esta experiencia fue la introducción de un componente estratégico: en determinadas sesiones, un miembro del grupo asumía un rol específico, introduciendo respuestas erróneas en las tareas. Esta dinámica implicaba que el resto del grupo debía prestar especial atención al desarrollo de la actividad, detectar posibles incoherencias y debatir en equipo. Este componente lúdico, con implicaciones cognitivas y sociales, favoreció el pensamiento crítico, la autorregulación emocional y la capacidad de análisis grupal.

Una vez finalizado el periodo de aplicación de la intervención gamificada, se administró el postest (CCLE) a toda la muestra y se recogieron los datos para su posterior análisis estadístico. Este análisis tuvo como propósito dar respuesta a las preguntas de investigación y valorar el cumplimiento de los objetivos planteados. Durante este proceso, se implementó una medida de equidad educativa: el grupo de control participó posteriormente en la experiencia gamificada, garantizando así la igualdad de oportunidades en el acceso a metodologías activas.

La investigación fue posteriormente validada por el Comité de Ética en Investigación Social de la Universidad de Castilla-La Mancha, lo que garantizó su adecuación ética y metodológica.

\section{Análisis de datos y resultados}
Para el análisis estadístico de los datos obtenidos se utilizó el \textit{software} SPSS (Statistical Package for the Social Sciences). En primer lugar, se verificó la distribución no normal de los datos mediante la prueba de Kolmogorov-Smirnov (p < 0,05). Posteriormente, se empleó la prueba Mann-Whitney para muestras independientes con el fin de identificar posibles diferencias significativas en función del grupo y del género. Asimismo, se calcularon las medias (M) y desviaciones estándar (DE) como parámetros descriptivos. Se estableció un nivel de significación estadística de p < 0,05.

\subsection{Diferencias de género en el rendimiento académico en competencia lingüística}
Se analizaron posibles diferencias en el rendimiento académico en competencia lingüística en función del género, tanto en el pretest como en el postest. Tal como se observa en la Tabla \ref{tab-2}, los resultados de la prueba de Mann-Whitney no evidenciaron diferencias estadísticamente significativas en ninguno de los momentos de medición (p > 0,05), lo que sugiere un rendimiento académico homogéneo entre niños y niñas.

%--- codigo da tabela 2 ---%
\begin{table}[htbp]
\centering
\begin{threeparttable}
\caption{Prueba Mann-Whitney: diferencias de género (total de la muestra).}\label{tab-2}
\begin{tabular}{p{3cm}p{2cm}p{2cm}p{2cm}}
\toprule
CCLE & U & Z & Sig. \\
\midrule
Pretest & 1112 & -0,293 & 0,769 \\
Postest & 1048 & -0,764 & 0,445 \\
\bottomrule
\end{tabular}
\source{Elaboración propia.}
\end{threeparttable}
\end{table}

\subsection{Diferencias de género por grupo}
Con el fin de explorar posibles diferencias de género específicas dentro de cada grupo, se aplicó la prueba no paramétrica de Mann-Whitney en los grupos de control y experimental. Tal como se observa en la Tabla \ref{tab-3}, no hubo diferencias estadísticamente significativas entre niños y niñas en el pretest (p > 0,05), lo que indica una homogeneidad académica inicial. No obstante, en el postest se observaron diferencias significativas por género dentro del grupo de control (p = 0,020), mientras que en el grupo experimental no hubo diferencias significativas (p = 0,321), lo que refuerza el posible ``efecto nivelador'' de la gamificación.

%--- codigo da tabela 3 ---%
\begin{table}[htbp]
\centering
\begin{threeparttable}
\caption{Prueba Mann-Whitney: diferencias de género según el grupo.}\label{tab-3}
\begin{tabular}{p{1.5cm}p{2cm}p{1cm}p{2.5cm}p{1cm}}
\toprule
CCLE & Grupo & U & Error estándar & Sig. \\
\midrule
Pretest & Control & 312 & 48,26 & 0,619 \\
& Experimental & 288 & 48,26 & 1,000 \\

\addlinespace

Postest & Control & 400 & 48,26 & 0,020 \\
& Experimental & 240 & 48,34 & 0,321 \\
\bottomrule
\end{tabular}
\source{Elaboración propia.}
\end{threeparttable}
\end{table}

\subsection{Comparación entre metodologías según el género}
Se compararon los resultados entre el grupo de control y el grupo experimental para cada género, antes y después de la intervención. En el pretest, hubo diferencias significativas ni para niños ni para niñas (p > 0,05), lo que sugiere condiciones equivalentes iniciales. En el postest, las niñas del grupo experimental mejoraron su rendimiento respecto del grupo de control, pero las diferencias no alcanzaron significación estadística (p = 0,319). Sin embargo, los niños del grupo experimental obtuvieron un rendimiento significativamente superior (U = 496; p < 0,001) frente a los del grupo de control (ver Tabla \ref{tab-4}). Esta mejora no generó diferencias entre niños y niñas dentro del grupo gamificado (p = 0,321).

%--- codigo da tabela 4 ---%
\begin{table}[h!]
\centering
\begin{threeparttable}
\caption{Prueba Mann-Whitney: diferencias metodológicas por género.}\label{tab-4}
\begin{tabular}{p{1.3cm}p{1.5cm}p{1cm}p{2.5cm}p{1.5cm}}
\toprule
Género & CCLE & U & Error estándar & Sig. \\
\midrule
Niños & Pretest & 352 & 48,26 & 0,185 \\
& Postest & 496 & 48 & < 0,001 \\

\addlinespace

Niñas & Pretest & 288 & 48,34 & 1,000 \\
& Postest & 336 & 48,17 & 0,319 \\
\bottomrule
\end{tabular}
\source{Elaboración propia.}
\end{threeparttable}
\end{table}

\subsection{Estadísticos descriptivos por grupo y género}
La Tabla \ref{tab-5} presenta los estadísticos descriptivos (media y desviación estándar) del rendimiento académico en competencia lingüística en inglés (CCLE), desglosados por género, grupo y momento de evaluación. Se destaca el rendimiento de los niños en el grupo experimental en el postest (M = 24,88; DE = 3,83), superior al de sus homólogos del grupo de control (M = 18,71; DE = 4,49), lo que refleja cuantitativamente el efecto significativo de la intervención didáctica. En el caso de las niñas, si bien también se observa una mejora en el grupo experimental, la diferencia entre grupos no fue significativa.

En cuanto a los valores del pretest, se constata que niños y niñas presentaban puntuaciones similares al inicio del estudio. En el grupo de control, la media fue ligeramente superior en las niñas (M = 12,96) frente a los niños (M = 11,17); patrón que se repite en el grupo experimental (niñas: M = 12,71; niños: M = 12,42). Estas diferencias iniciales sugieren una leve ventaja femenina en el rendimiento lingüístico. En el postest, las niñas del grupo de control mantuvieron un rendimiento superior al de los niños (M = 21,46 frente a M = 18,71), diferencia que resultó estadísticamente significativa según se discutió en el apartado anterior. En el grupo experimental, los niños alcanzan una media superior (M = 24,88) a la de sus compañeras (M = 23,33).

%--- codigo da tabela 5 ---%
\begin{table}[h!]
\centering
\begin{threeparttable}
\caption{Estadísticos descriptivos: media (M) y desviación estándar (DE).}\label{tab-5}
\begin{tabular}{p{1.3cm}p{1.5cm}p{2.5cm}p{1cm}p{1cm}p{0.8cm}}
\toprule
Género & CCLE & Grupo & N & M & DE \\
\midrule
Niños & Pretest & Control & 24 & 11,17 & 4,01 \\
& & Experimental & 24 & 12,42 & 2,47 \\
& Postest & Control & 24 & 18,71 & 4,49 \\
& & Experimental & 24 & 24,88 & 3,83 \\

\addlinespace

Niñas & Pretest & Control & 24 & 12,96 & 7,75 \\
& & Experimental & 24 & 12,71 & 7,94 \\
& Postest & Control & 24 & 21,46 & 6,10 \\
& & Experimental & 24 & 23,33 & 4,61 \\
\bottomrule
\end{tabular}
\source{Elaboración propia.}
\end{threeparttable}
\end{table}

\section{Discusión}

\subsection{Rendimiento académico en lengua inglesa y género: resultados iniciales}
Los resultados del pretest, tras la aplicación de la prueba Mann-Whitney, muestran que no hubo diferencias significativas entre niños y niñas en el rendimiento inicial en competencia lingüística (ver Tabla \ref{tab-2}), lo que apunta a una homogeneidad académica entre niños y niñas. Al desagregar por grupo, tampoco se detectaron diferencias significativas en el pretest, ni en el grupo de control (U = 312; p = 0,619) ni en el experimental (U = 288; p = 1,000). Este hallazgo permite atribuir cualquier efecto posterior observado al tipo de metodología implementada y no a desigualdades iniciales entre géneros.

Los análisis preliminares confirman que la muestra partía de condiciones equivalentes, tanto en el conjunto global como en cada grupo. Esta homogeneidad se alinea con estudios previos que afirman que, en edades tempranas, las diferencias de género en el rendimiento lingüístico tienden a ser mínimas o inexistentes cuando se controlan factores contextuales, como la influencia cultural o sociopsicológica \cite{saville2012}.

Sin embargo, las puntuaciones de las niñas son ligeramente superiores a las de los niños en el pretest: en el grupo de control, las niñas obtuvieron una puntuación de 12,96, frente a la de los niños, que fue casi de 2 puntos menos (M = 11,17); patrón que se repite en el grupo experimental con un aumento ligeramente mayor (niñas: M = 12,71; niños: M = 12,42). Estas diferencias iniciales coinciden con investigaciones previas que señalan una ligera ventaja femenina en el rendimiento lingüístico \cite{barrios2015, glowka2014}.

\subsection{Efectos de la gamificación en el rendimiento en lengua inglesa según el género}
El análisis de las medias y desviaciones estándar refuerza el hallazgo de que la gamificación impacta de forma diferencial según el género. En el caso de los niños, los valores descriptivos en el pretest fueron comparables: control (M = 11,17; DE = 4,01) y experimental (M = 12,42; DE = 2,47), sin significancia (U = 352; p = 0,185). Tras la intervención, la media de los varones del grupo experimental (M = 24,88; DE = 3,83) superó notablemente a la de sus homólogos en el grupo de control (M = 18,71; DE = 4,49), resultando en diferencia significativa (U = 496; p < 0,001). Esta mejora sugiere que la gamificación, centrada en narrativa interactiva y refuerzo positivo, puede potenciar la implicación y el rendimiento de estudiantes que tradicionalmente han mostrado menor motivación hacia tareas de tipo lingüístico \cite{prieto2018, dicheva2018, bartholomew2018}. \textcite{fiestas2023} destacan que el efecto de las metodologías activas puede ser más visible en alumnado con motivación inicial baja, como es frecuente en los niños que aprenden una segunda lengua.

Por su parte, las niñas mostraron un avance moderado: en el pretest, el grupo de control (M = 12,96; DE = 7,75) y el experimental (M = 12,71; DE = 7,94) resultaron prácticamente idénticos (U = 288; p = 1,000), y en el postest la media creció en ambos grupos (control: M = 21,46; DE = 6,10; experimental: M = 23,33; DE = 4,61) sin que la diferencia fuera significativa (U = 336; p = 0,319). Estos resultados pueden deberse a que las niñas tienden a mostrar una motivación más intrínseca hacia la lengua extranjera, independientemente del enfoque didáctico \cite{fernandezbarrionuevo2022, abdullahi2015}. 

En conjunto, estos resultados ponen de relieve un ``efecto nivelador'' de la gamificación. Al favorecer mejoras significativas en el rendimiento académico del alumnado masculino, esta metodología equiparó sus resultados a los de sus compañeras, quienes ya tenían alto desempeño y siguieron mejorando con la metodología ordinaria. De hecho, no se observaron diferencias significativas entre géneros en el postest, lo que indica que la gamificación redujo disparidades iniciales y alcanzó un equilibrio en los logros académicos. Este efecto ya ha sido documentado en pedagogías andamiadas que equilibran el rendimiento del alumnado \cite{nieto2019, nieto2026, nieto-casanova2026}.

\subsection{Diferencias de género tras la intervención}
Tras la intervención, el contraste global por género en el postest confirma que no persisten disparidades significativas al considerar toda la muestra (U = 1048; p = 0,445). Al evaluar cada metodología por separado, surge un matiz: en el grupo de control, las niñas alcanzaron una media de 21,46 (DE = 6,10) frente a 18,71 (DE = 4,49) de los niños, diferencia que sí resultó significativa (U = 400; p = 0,020). En cambio, en el grupo experimental, chicos (M = 24,88; DE = 3,83) y chicas (M = 23,33; DE = 4,61) no mostraron diferencias estadísticamente relevantes (U = 240; p = 0,321).

Estos datos evidencian que la metodología tradicional tiende a reproducir o acentuar brechas de género, mientras que la gamificación las neutraliza, corroborando que las prácticas gamificadas favorecen la equidad de resultados \cite{zahedi2021}. Esta conclusión es coherente con \textcite{navarro2021}, quienes argumentan que la gamificación, cuando se diseña atendiendo a los intereses y motivaciones del alumnado, genera escenarios de aprendizaje más democráticos y motivadores. Al combinar elementos lúdicos y retos cognitivos, esta metodología implica de forma significativa a estudiantes con estilos de aprendizaje diversos, favoreciendo especialmente a quienes requieren un componente emocional o motivacional más fuerte para comprometerse con tareas lingüísticas, como los niños \cite{navarro2021, williams2013}.

\section{Conclusiones, implicaciones pedagógicas y futuras líneas de investigación}
El estudio evaluó el impacto de una metodología gamificada con apoyo digital sobre el rendimiento en lengua inglesa de alumnado de 2.º de EP, con especial atención a diferencias de género, mediante un diseño cuasi-experimental pretest-postest con dos grupos (control y experimental) y una muestra de 96 estudiantes (50\% niños, 50\% niñas), aplicando el \textit{Cuestionario de Competencias Lingüísticas Específicas (CCLE)}.

Independientemente de la metodología, no se detectaron diferencias de género en el pretest, lo que valida que los efectos observados en el postest se deben a la intervención. La gamificación mejoró significativamente el rendimiento de los estudiantes varones, equiparándolo al de las niñas, cuya trayectoria no se vio afectada por su aptitud inicial. En este sentido, los resultados obtenidos muestran que la gamificación, diseñada según los intereses del alumnado, incide positivamente en el rendimiento académico, especialmente en el alumnado masculino.

Estos hallazgos respaldan investigaciones previas que evidencian el potencial de la gamificación para aumentar la motivación y la implicación del alumnado, siendo particularmente eficaz en perfiles con menor motivación intrínseca hacia el aprendizaje de lenguas extranjeras, como es común en parte del alumnado masculino. El ``efecto nivelador'' se confirma al reducir las brechas de género observadas bajo metodologías convencionales, consolidando la gamificación como andamiaje pedagógico equitativo.

Desde una perspectiva pedagógica, los resultados invitan a reflexionar sobre el diseño de entornos de aprendizaje inclusivos, flexibles y adaptativos, que integren mecánicas de gamificación de forma planificada y fundamentada. A diferencia de enfoques más rígidos, la gamificación permite atender diferentes estilos de aprendizaje, minimizando sesgos de género y favoreciendo la equidad en el aula.

Entre las implicaciones prácticas, se destaca la necesidad de formación docente específica orientada al diseño e implementación de experiencias gamificadas alineadas con los objetivos curriculares y las competencias lingüísticas. Asimismo, se recomienda usar plataformas digitales accesibles, como Genially o ClassDojo, que permitan construir narrativas didácticas inmersivas y ofrecer sistemas de retroalimentación motivadora, así como desarrollar sistemas de evaluación coherentes con este enfoque metodológico que valoren el rendimiento lingüístico, el grado de implicación y la colaboración.

Para futuras investigaciones, se sugiere analizar los efectos longitudinales de la gamificación sobre el rendimiento y la motivación del alumnado en el aprendizaje de lenguas extranjeras. A ello se suma la necesidad de explorar su interacción con variables psicoeducativas como la autoeficacia, la orientación motivacional o la ansiedad comunicativa. Asimismo, se reconoce que el tamaño de la muestra limita la representatividad de los datos. Por ello, futuras investigaciones podrían ampliar la muestra, incorporando contextos educativos diversos y un mayor número de participantes. Finalmente, se propone incorporar enfoques metodológicos cualitativos o mixtos que permitan recoger las percepciones del alumnado y del profesorado sobre su experiencia con la gamificación, aportando así una visión más integral del fenómeno estudiado.

En síntesis, esta investigación aporta evidencia empírica que refuerza el valor pedagógico de la gamificación como herramienta eficaz para promover el rendimiento académico en lengua inglesa en contextos escolares, especialmente cuando se busca maximizar la implicación del alumnado y fomentar la equidad metodológica. El desafío educativo actual no radica únicamente en enseñar contenidos lingüísticos, sino en crear espacios de aprendizaje significativos, motivadores y sostenibles para todos los perfiles del alumnado. Esta constatación invita a reflexionar sobre el \textit{statu quo} de la enseñanza del inglés como lengua extranjera en la EP y plantea un interrogante de calado pedagógico: ¿En qué medida la integración de la gamificación en la enseñanza del inglés puede transformar las prácticas educativas convencionales y abrir nuevas vías de equidad y eficacia en el aprendizaje lingüístico?

\section*{Agradecimientos}
Esta investigación forma parte del proyecto \textit{Mejora de los procesos de enseñanza de lenguas: protocolos de actuación y experimentación para la enseñanza bilingüe familiar (PLF), escolar (AICLE) y universitaria (EMI) y para la innovación didáctica} (referencia 2022-GRIN-34455), financiado por la UCLM y el FEDER. Asimismo, se extiende el agradecimiento al Departamento de Filología Moderna de la UCLM y a los centros de Educación Primaria participantes, incluyendo equipos directivos, profesorado, familias y alumnado.


\printbibliography\label{sec-bib}
% if the text is not in Portuguese, it might be necessary to use the code below instead to print the correct ABNT abbreviations [s.n.], [s.l.]
%\begin{portuguese}
%\printbibliography[title={Bibliography}]
%\end{portuguese}


%full list: conceptualization,datacuration,formalanalysis,funding,investigation,methodology,projadm,resources,software,supervision,validation,visualization,writing,review

\begin{dataavailability}
\txtdataavailability{databody} % options: dataavailable, dataonly, databody, datanotav, nodata
\end{dataavailability}





\end{document}

